\section{Bộ nhớ}

\subsection{Câu hỏi nghiên cứu}

\textit{Memory có nên chỉ là database hay nên trở thành một learning system?}

\subsection{Key Findings}

\begin{enumerate}
    \item \textbf{Hybrid approach đạt 91\% F1} --- kết hợp vector retrieval + LLM reranking
    \item \textbf{Memory decay là exponential} --- 1.3pp/day giảm accuracy
    \item \textbf{Generalization gap lớn} --- MemOS: 90\% trên LoCoMo $\to$ 55\% trên LifeBench ($-34.78$pp)
    \item \textbf{Forgetting mechanisms cần thiết} --- FactConsolidation chỉ đạt 54\% single-hop accuracy
\end{enumerate}

\begin{figure}[H]
\centering
\includegraphics[width=0.95\textwidth]{figures/fig02-memory-comparison.pdf}
\caption{So sánh kiến trúc Memory --- Hybrid (vector + graph + LLM rerank) đạt accuracy cao nhất (91\% F1) với latency chấp nhận được}
\label{fig:memory-comparison}
\end{figure}

\subsection{Khuyến nghị}

Memory nên là \textbf{learning system} với các components:
\begin{itemize}
    \item \textbf{Vector DB} (ChromaDB/Pinecone) cho episodic memories
    \item \textbf{Graph DB} (Neo4j) cho semantic relationships
    \item \textbf{Consolidation pipeline} cho periodic summary generation
    \item \textbf{Forgetting curve} cho automatic pruning
\end{itemize}

% =============================================================================
% 9. RELATIONSHIP